\newpage
\appendix
\section*{Appendix}
\addcontentsline{toc}{section}{Appendix}
\renewcommand{\thesubsection}{A.\arabic{subsection}}
\renewcommand{\thesubsubsection}{A.\arabic{subsection}.\arabic{subsubsection}}

\section{Prompting templates}
Our prompt is composed of a graph definition, function descriptions, and demonstrations. The complete prompt is shown below, where \texttt{\{graph\_definition\}}, \texttt{\{function\_definitions\}}, and \texttt{\{demonstrations\}} correspond to each respective section in the prompt.


\subsection{Full prompt}
Solve a question answering task with interleaving Thought, Interaction with Graph, and Feedback from Graph steps. 

In the Thought step, you should reflect on what further information is needed to answer the question. In the Interaction step, you interact with the graph using one of the following four functions:

\texttt{\{function\_definitions\}}

Here are some examples. You should follow the structure of the examples in your output. 

\texttt{\{demonstrations\}}




\vspace{0.2cm}
Graph Definition: \texttt{\{graph\_definition\}} \\
Question: \texttt{\{question\}} 

\vspace{0.2cm}
Please answer by providing node main features (e.g., names) rather than node IDs. YOU MUST follow the structure of the examples.



\subsection{Graph definitions}
The graph definition describes the structure of each knowledge graph. This includes information like how many types of nodes are in the graph and their associated features.  

\subsubsection{Maple graph description}
There are three types of nodes in the graph: paper, author and venue.\\
Paper nodes have features: title, abstract, year and label. Author nodes have features: name. Venue nodes have features: name.\\
Paper nodes are linked to author nodes, venue nodes, reference nodes and cited\_by nodes. Author nodes are linked to paper nodes. Venue nodes are linked to paper nodes.

\subsubsection{Biomedical graph description}
There are eleven types of nodes in the graph: Anatomy, Biological Process, Cellular Component, Compound, Disease, Gene, Molecular Function, Pathway, Pharmacologic Class, Side Effect, Symptom.\\
Each node has name feature.\\
There are these types of edges: Anatomy-downregulates-Gene, Anatomy-expresses-Gene, Anatomy-upregulates-Gene, Compound-binds-Gene, Compound-causes-Side Effect, Compound-downregulates-Gene, Compound-palliates-Disease, Compound-resembles-Compound, Compound-treats-Disease, Compound-upregulates-Gene, Disease-associates-Gene, Disease-downregulates-Gene, Disease-localizes-Anatomy, Disease-presents-Symptom, Disease-resembles-Disease, Disease-upregulates-Gene, Gene-covaries-Gene, Gene-interacts-Gene, Gene-participates-Biological Process, Gene-participates-Cellular Component, Gene-participates-Molecular Function, Gene-participates-Pathway, Gene-regulates-Gene, Pharmacologic Class-includes-Compound.

\subsubsection{Legal graph description}
There are four types of nodes in the graph: opinion, opinion\_cluster, docket, and court.\\
Opinion nodes have features: plain\_text. Opinion\_cluster nodes have features: syllabus, judges, case\_name, attorneys. Docket nodes have features: pacer\_case\_id, case\_name. Court nodes have features: full\_name, start\_date, end\_date, citation\_string.\\
Opinion nodes are linked to their reference nodes and cited\_by nodes, as well as their opinion\_cluster nodes. Opinion\_cluster nodes are linked to opinion nodes and docket nodes. Docket nodes are linked to opinion\_cluster nodes and court nodes. Court nodes are linked to docket nodes.

\subsubsection{Amazon graph description}
There are two types of nodes in the graph: item and brand.\\
Item nodes have features: title, description, price, img, category. Brand nodes have features: name.\\
Item nodes are linked to their brand nodes, also\_viewed\_item nodes, buy\_after\_viewing\_item nodes, also\_bought\_item nodes, bought\_together\_item nodes. Brand nodes are linked to their item nodes.

\subsubsection{Goodreads graph description}
There are four types of nodes in the graph: book, author, publisher, and series.\\
Book nodes have features: country\_code, language\_code, is\_ebook, title, description, format, num\_pages, publication\_year, url, popular\_shelves, and genres. Author nodes have features: name. Publisher nodes have features: name. Series nodes have features: title and description.\\
Book nodes are linked to their author nodes, publisher nodes, series nodes and similar\_books nodes. Author nodes are linked to their book nodes. Publisher nodes are linked to their book nodes. Series nodes are linked to their book nodes.

\subsubsection{DBLP graph description}
There are three types of nodes in the graph: paper, author and venue.\\
Paper nodes have features: title, abstract, keywords, lang, and year. Author nodes have features: name and organization. Venue nodes have features: name.\\
Paper nodes are linked to their author nodes, venue nodes, reference nodes (the papers this paper cite) and cited\_by nodes (other papers which cite this paper). Author nodes are linked to their paper nodes. Venue nodes are linked to their paper nodes.

\subsection{Function Definitions}

You are a helpful AI bot who can assist in answering questions with help of a knowledge graph. You can answer questions by interacting with the graph using four functions:

(1) RetrieveNode[keyword], which retrieves the related node from the graph according to the corresponding query.

(2) NodeFeature[Node, feature], which returns the detailed attribute information of Node regarding the given "feature" key.

(3) NodeDegree[Node, neighbor\_type], which calculates the number of "neighbor\_type" neighbors of the node Node in the graph.

(4) NeighbourCheck[Node, neighbor\_type], which lists the "neighbor\_type" neighbours of the node Node in the graph and returns them.

You answer these questions with an interleaving Thought, Action taken on the graph using one of the above functions, and an Observation that is returned from the graph. You may take as many steps as necessary to answer the question.

In the Thought step you should come up with a plan to answer the question, this can include what information is still needed or a plan to solving the question. 
In the Action step you should interact with the graph using one of the four functions based on the previous Thought. 
In the Observation step you should provide the information returned from the graph based on the previous Action.

If nothing is returned from the graph, you should try another strategy. e.g. Observation 1: the reference neighbors of 123456 are: []. You should move on to another strategy. 


\subsection{Demonstrations}
BEGINNING OF EXAMPLES

\textbf{Example 1} \par
\medskip % Adds a small vertical space after the "Example N" title

\textbf{Definition of the graph:}
There are three types of nodes in the graph: paper, author and venue. \par
Paper nodes have features: title, abstract, year and label. Author nodes have features: name. Venue nodes have features: name. \par
Paper nodes are linked to their author nodes, venue nodes, reference paper nodes and cited\_by paper nodes. Author nodes are linked to their paper nodes. Venue nodes are linked to their paper nodes. \par

\textbf{Question:}
When was the paper Strongly Interacting Higgs Sector in the Minimal Standard Model published? \par

\textbf{Thought 1:} The question is asking some basic information about a paper. I should first find the paper node related to the paper Strongly Interacting Higgs Sector in the Minimal Standard Model. \par
\textbf{Action 1:} \texttt{RetrieveNode[Strongly Interacting Higgs Sector in the Minimal Standard Model]} \par
\textbf{Observation 1:} The ID of this node is \texttt{3101448248}. \par

\textbf{Thought 2:} The question is asking for the publication year of the paper. I need to check the node feature (year) from the graph. \par
\textbf{Action 2:} \texttt{NodeFeature[3101448248, year]} \par
\textbf{Observation 2:} 1993 \par

\textbf{Thought 3:} The publication year of the paper is 1993. \par
\textbf{Action 3:} \texttt{Finish[1993]} \par

\bigskip % Adds more vertical space between examples

\textbf{Example 2} \par
\medskip

\textbf{Definition of the graph:}
There are three types of nodes in the graph: paper, author and venue. \par
Paper nodes have features: title, abstract, year and label. Author nodes have features: name. Venue nodes have features: name. \par
Paper nodes are linked to their author nodes, venue nodes, reference paper nodes and cited\_by paper nodes. Author nodes are linked to their paper nodes. Venue nodes are linked to their paper nodes. \par

\textbf{Question:}
How many authors do the paper Mass Accretion Rates in Self-Regulated Disks of T Tauri Stars have? \par

\textbf{Thought 1:} The question is asking for the number of authors of a paper. I should first find the paper node related to the paper Mass Accretion Rates in Self-Regulated Disks of T Tauri Stars in the graph. \par
\textbf{Action 1:} \texttt{RetrieveNode[Mass Accretion Rates in Self-Regulated Disks of T Tauri Stars]} \par
\textbf{Observation 1:} The ID of this node is \texttt{2090642949}. \par

\textbf{Thought 2:} The question is asking for the number of authors of the paper. I should check the node's author neighbor degree from the graph. \par
\textbf{Action 2:} \texttt{NodeDegree[2090642949, author]} \par
\textbf{Observation 2:} 5 \par

\textbf{Thought 3:} The number of authors of the paper is 5. \par
\textbf{Action 3:} \texttt{Finish[5]} \par

\bigskip

\textbf{Example 3} \par
\medskip

\textbf{Definition of the graph:}
There are three types of nodes in the graph: paper, author and venue. \par
Paper nodes have features: title, abstract, year and label. Author nodes have features: name. Venue nodes have features: name. \par
Paper nodes are linked to their author nodes, venue nodes, reference paper nodes and cited\_by paper nodes. Author nodes are linked to their paper nodes. Venue nodes are linked to their paper nodes. \par

\textbf{Question:}
What was the publish venue of the paper Mass Accretion Rates in Self-Regulated Disks of T Tauri Stars? \par

\textbf{Thought 1:} The question is asking information about the node related to the paper Mass Accretion Rates in Self-Regulated Disks of T Tauri Stars. I should first find the paper node related to the paper in the graph. \par
\textbf{Action 1:} \texttt{RetrieveNode[Mass Accretion Rates in Self-Regulated Disks of T Tauri Stars]} \par
\textbf{Observation 1:} The ID of this node is \texttt{2090642949}. \par

\textbf{Thought 2:} The question is asking for the publish venue of the paper. I should check the node's venue neighbor feature from the graph. \par
\textbf{Action 2:} \texttt{NodeFeature[2090642949, venue]} \par
\textbf{Observation 2:} the astrophysical journal, the atmosphere journal \par

\textbf{Thought 3:} The publish venue of the paper is the astrophysical journal, the atmosphere journal. \par
\textbf{Action 3:} \texttt{Finish[the astrophysical journal, the atmosphere journal]} \par

\bigskip

\textbf{Example 4} \par
\medskip

\textbf{Definition of the graph:}
There are three types of nodes in the graph: paper, author and venue. \par
Paper nodes have features: title, abstract, year and label. Author nodes have features: name. Venue nodes have features: name. \par
Paper nodes are linked to their author nodes, venue nodes, reference paper nodes and cited\_by paper nodes. Author nodes are linked to their paper nodes. Venue nodes are linked to their paper nodes. \par

\textbf{Question:}
Who wrote the paper titled ``Strongly Interacting Higgs Sector in the Minimal Standard Model''? \par

\textbf{Thought 1:} The question is asking for the author(s) of a paper. I should first find the paper node related to the paper titled ``Strongly Interacting Higgs Sector in the Minimal Standard Model'' in the graph. \par
\textbf{Action 1:} \texttt{RetrieveNode[Strongly Interacting Higgs Sector in the Minimal Standard Model]} \par
\textbf{Observation 1:} The ID of this node is \texttt{3101448248}. \par

\textbf{Thought 2:} The question is asking for the author(s) of the paper. I should check the node's author neighbor degree from the graph. \par
\textbf{Action 2:} \texttt{NeighbourCheck[3101448248, author]} \par
\textbf{Observation 2:} The author neighbors of \texttt{3101448248} are: [\texttt{'98581309'}, \texttt{'2440536741'}, \texttt{'2345070782'}, \texttt{'2113669516'}, \texttt{'1968129287'}, \texttt{'2112021749'}, \texttt{'2165059783'}, \texttt{'2790497238'}, \texttt{'2035459489'}, \texttt{'2166580174'}]. \par

\textbf{Thought 3:} The author(s) of the paper are: [\texttt{'98581309'}, \texttt{'2440536741'}, \texttt{'2345070782'}, \texttt{'2113669516'}, \texttt{'1968129287'}, \texttt{'2112021749'}]. I should check each of the author nodes to get the author names. \par
\textbf{Action 3:} \texttt{NodeFeature[98581309, name]} \par
\textbf{Observation 3:} Author 1: John F. Gunion \par

\textbf{Thought 4:} The author(s) of the paper are: [\texttt{'98581309'}, \texttt{'2440536741'}, \texttt{'2345070782'}, \texttt{'2113669516'}, \texttt{'1968129287'}, \texttt{'2112021749'}]. I should check each of the author nodes to get the author names. I have already checked the first author. I should check the next author. \par
\textbf{Action 4:} \texttt{NodeFeature[2440536741, name]} \par
\textbf{Observation 4:} Author 2: Howard E. Haber \par

\textbf{Thought 5:} The author(s) of the paper are: [\texttt{'98581309'}, \texttt{'2440536741'}, \texttt{'2345070782'}, \texttt{'2113669516'}, \texttt{'1968129287'}, \texttt{'2112021749'}]. I should check each of the author nodes to get the author names. I have already checked the first two authors. I should check the next author. \par
\textbf{Action 5:} \texttt{NodeFeature[2345070782, name]} \par
\textbf{Observation 5:} Author 3: Gordon Kane \par

\textbf{Thought 6:} The author(s) of the paper are: [\texttt{'98581309'}, \texttt{'2440536741'}, \texttt{'2345070782'}, \texttt{'2113669516'}, \texttt{'1968129287'}, \texttt{'2112021749'}]. I should check each of the author nodes to get the author names. I have already checked the first three authors. I should check the next author. \par
\textbf{Action 6:} \texttt{NodeFeature[2113669516, name]} \par
\textbf{Observation 6:} Author 4: Sally Dawson \par

\textbf{Thought 7:} The author(s) of the paper are: [\texttt{'98581309'}, \texttt{'2440536741'}, \texttt{'2345070782'}, \texttt{'2113669516'}, \texttt{'1968129287'}, \texttt{'2112021749'}]. I should check each of the author nodes to get the author names. I have already checked the first four authors. I should check the next author. \par
\textbf{Action 7:} \texttt{NodeFeature[1968129287, name]} \par
\textbf{Observation 7:} Author 5: Michael S. Chanowitz \par

\textbf{Thought 8:} The author(s) of the paper are: [\texttt{'98581309'}, \texttt{'2440536741'}, \texttt{'2345070782'}, \texttt{'2113669516'}, \texttt{'1968129287'}, \texttt{'2112021749'}]. I should check each of the author nodes to get the author names. I have already checked the first five authors. I should check the next author. \par
\textbf{Action 8:} \texttt{NodeFeature[2112021749, name]} \par
\textbf{Observation 8:} Author 6: Lian-Tao Wang \par

\textbf{Thought 9:} The author(s) of the paper are: [\texttt{'98581309'}, \texttt{'2440536741'}, \texttt{'2345070782'}, \texttt{'2113669516'}, \texttt{'1968129287'}, \texttt{'2112021749'}]. I have checked all the authors and am finished. \par
\textbf{Action 9:} \texttt{Finish[John F. Gunion, Howard E. Haber, Gordon Kane, Sally Dawson, Michael S. Chanowitz, Lian-Tao Wang]} \par

\bigskip

\textbf{Example 5} \par
\medskip

\textbf{Definition of the graph:}
There are three types of nodes in the graph: paper, author and venue. \par
Paper nodes have features: title, abstract, year and label. Author nodes have features: name. Venue nodes have features: name. \par
Paper nodes are linked to their author nodes, venue nodes, reference paper nodes and cited\_by paper nodes. Author nodes are linked to their paper nodes. Venue nodes are linked to their paper nodes. \par

\textbf{Question:}
Considering a reader of the paper Strongly Interacting Higgs Sector in the Minimal Standard Model, Select another paper that the reader would be interested in reading. \par

\textbf{Thought 1:} The question is asking for the papers that a reader of the paper Strongly Interacting Higgs Sector in the Minimal Standard Model would be interested in reading. I should first find the paper node related to the paper Strongly Interacting Higgs Sector in the Minimal Standard Model in the graph. \par
\textbf{Action 1:} \texttt{RetrieveNode[Strongly Interacting Higgs Sector in the Minimal Standard Model]} \par
\textbf{Observation 1:} The ID of this node is \texttt{3101448248}. \par

\textbf{Thought 2:} I should also check the venue of the paper Strongly Interacting Higgs Sector in the Minimal Standard Model to provide the best paper. \par
\textbf{Action 2:} \texttt{NodeFeature[3101448248, venue]} \par
\textbf{Observation 2:} The venue of the paper is Physical Review D. \par

\textbf{Thought 3:} The question is asking for the papers that a reader of the paper would be interested in reading. I should check the node's reference paper neighbors from the graph. \par
\textbf{Action 3:} \texttt{NeighbourCheck[3101448248, reference]} \par
\textbf{Observation 3:} The reference paper neighbors of \texttt{3101448248} are: [\texttt{'1988913734'}, \texttt{'1988913735'}, \texttt{'1988913736'}]. \par

\textbf{Thought 4:} The papers that a reader of the paper would be interested in reading are: [\texttt{'1988913734'}, \texttt{'1988913735'}, \texttt{'1988913736'}]. I should check the first paper. \par
\textbf{Action 4:} \texttt{NodeFeature[1988913734, title]} \par
\textbf{Observation 4:} Paper 1: Advances in Quantum Computing Algorithms \par

\textbf{Thought 5:} The papers that a reader of the paper would be interested in reading are: [\texttt{'1988913735'}, \texttt{'1988913736'}]. I should check the second paper. \par
\textbf{Action 5:} \texttt{NodeFeature[1988913735, title]} \par
\textbf{Observation 5:} Paper 2: Machine Learning Techniques for Big Data Analysis \par

\textbf{Thought 6:} The papers that a reader of the paper would be interested in reading are: [\texttt{'1988913734'}, \texttt{'1988913735'}, \texttt{'1988913736'}]. I should check the third paper. \par
\textbf{Action 6:} \texttt{NodeFeature[1988913736, title]} \par
\textbf{Observation 6:} Paper 3: The Higgs Sector and Hidden Symmetry \par

\textbf{Thought 7:} The venue of the papers Strongly Interacting Higgs Sector in the Minimal Standard Model and The Higgs Sector and Hidden Symmetry are Physical Review D. which is the same. The question asked for one paper so I will select the best paper. \par
\textbf{Action:} \texttt{Finish[The Higgs Sector and Hidden Symmetry]} \par

END OF EXAMPLES


\section{Hyperparameteres}
\label{sec:appendix_hyperparameters}

This section details the hyperparameters and implementation specifics used in our experiments.

For retrieval and indexing, we used MPNet-v2 as the retriever model for semantic retrieval of nodes from the knowledge graphs. Node embeddings were indexed using FAISS \cite{johnson2019billion} for efficient similarity search.

For LLM text generation, the following parameters were consistently applied: a Temperature of 0.7 and a Top-p (nucleus) sampling threshold of 0.9. These settings were chosen to balance diversity and coherence in the generated responses.

As detailed in Section 4.3 and evaluated throughout Section 4, we varied two primary inference scaling parameters. The Number of reasoning steps (Sequential Scaling) typically varied between 10, 25, and 50 steps. The Number of sampled thought-action pairs for majority voting (Parallel Scaling) typically varied between 1 (no parallel scaling/GraphCoT equivalent), 4, 8, and 16 votes. Specific configurations are detailed in Tables 1, 2, and 3.

All experiments were conducted using Python 3.9 and the Huggingface Transformers library (version 4.51.3). All experiments were conducted on NVIDIA A100 and H100 GPUs.


